\section{Trace/Run Correspondence}

The central bridge theorem relates three objects for a regular expression \(r\)
and input word \(w\):
\begin{enumerate}
\item successful leaves of the backtracking tree \(\trace(r,w)\);
\item marked matches of the positioned expression \(\mathsf{label}(r)\);
\item accepting runs of the position automaton built from \(\mathsf{label}(r)\).
\end{enumerate}

\begin{theorem}[Successful trace leaves correspond to accepting runs]
\label{thm:trace-run}
For every regular expression \(r\) in the regular core and every word \(w\),
there is a bijection between \(\mathsf{SuccLeaves}(\trace(r,w))\) and the
accepting marked runs of the position automaton for \(\mathsf{label}(r)\) over
\(w\).  Consequently,
\[
  |\mathsf{SuccLeaves}(\trace(r,w))|
  =
  \amb_{\nfa_{\mathsf{pos}}(r)}(w).
\]
\end{theorem}

The proof decomposes into two parts.  First, a successful trace leaf induces a
marked word by recording the atom positions read along that branch.  The
well-formedness of the trace rules implies that consecutive positions respect
the followpos relation, so the marked word is accepted by the position
automaton.  Second, an accepting marked run determines a unique successful
branch in the full backtracking tree, because the tree semantics enumerates all
alternatives allowed by the syntax.  The two translations are inverse up to the
irrelevant administrative nodes of the trace tree.

Priority is intentionally not erased completely.  The bijection preserves the
order induced by backtracking choices.  Therefore the first accepting run in
trace order corresponds to the match returned by the JavaScript-style engine.
Ambiguity counts how many successful alternatives exist; priority chooses one of
them.
